\documentclass[10pt,a4paper]{amsart}
\usepackage[margin=19mm]{geometry}
\usepackage{amsmath,amssymb,mathtools}
\usepackage[scr=rsfs,cal=euler]{mathalfa}
\usepackage{xfrac}
\usepackage[unicode,hidelinks,bookmarks=false]{hyperref}
\newcommand{\ie}{i.e.\ }
\newcommand{\cf}{cf.\ }
\newcommand{\resp}{respectively\ }
\newcommand{\Subsection}{\subsection}
\providecommand{\nobreakdash}{\nobreak\hbox{-}}
\setlength{\emergencystretch}{2em}
\multlinegap0pt
\usepackage{bm}
\usepackage{amssymb}
\usepackage{mathtools}
\usepackage[shortlabels]{enumitem}
\newtheorem{lemma}{Lemma}
\title{The Gao-Zhuang conjecture for the Heisenberg group over $\mathbb{F}_p$}
\author{Yongke Qu, Guoqing Wang}
\date{Source: arXiv:2608.23319v2 (CC BY 4.0)}
\begin{document}
\maketitle

\noindent\emph{Source and licence.} The Gao-Zhuang conjecture for the Heisenberg group over $\mathbb{F}_p$. Version arXiv:2608.23319v2 (CC BY 4.0), distributed by arXiv under the Creative Commons Attribution 4.0 International licence, http://creativecommons.org/licenses/by/4.0/, as recorded in the arXiv licence metadata for this exact version and shown on the abstract page https://arxiv.org/abs/2608.23319v2. Full licence notice: \texttt{licenses/arxiv-2608.23319-CC-BY-4.0.txt}.\par\smallskip

\noindent\textit{Editorial scope.} Counted unit: the article's lemma lemma:relative (source0.tex lines 449--457). The article's complete original proof of it is delivered with it (source0.tex lines 459--560, one \texttt{proof} environment of 102 lines). No mathematics is written, repaired or re-derived: the statement and the proof are the article's own lines, in the article's own notation, and no retained line is substituted or re-flowed.  Retained uncounted as the source-local context the proof needs: lines 423--424.  Nothing was omitted from any retained span, so the package carries no omission record.  No retained line contains a citation, so the wrapper carries no bibliography and no bibliography entry.  No retained line contains a figure environment, an image-inclusion command or drawing code, so no image is delivered and the package declares no asset dependency.  Editorial formatting only, in addition: an AMS \texttt{amsart} preamble that restates the article's own declarations and macros used by the retained lines, namely the environments \texttt{lemma} on separate counters, together with the standard packages those lines need.  The retained environments print in this document as Lemma 1 in order of appearance, whereas the article prints the same items with the numbering of its own full document.  The complete native source of the article as distributed in the arXiv e-print for this exact version is bundled unchanged as \texttt{sources/208405/source0.tex}.
\begin{lemma}\label{lemma:multilinearreduction} Let $N\ge 1$, and let $p$ be a prime. Let $f(X_1,\ldots,X_N)\in \mathbb{F}_p[X_1,\ldots,X_N]$ be a multilinear polynomial. If $f(X_1,\ldots,X_N)$ vanishes at every point $x\in\{0,1\}^N$, then $f$ is the zero polynomial. In particular, if $g\in \mathbb{F}_p[X_1,\ldots,X_N]$ is a multilinear polynomial such that $f(x)=g(x)$ for every $x\in\{0,1\}^N$, then $f=g$.
\end{lemma}

\begin{lemma}\label{lemma:relative}
Let
$A=a_1\boldsymbol{\cdot}\ldots\boldsymbol{\cdot}a_N$
be a sequence over $V=C_p^2$ with $N<3p$, and let $H\le V$ be a
subgroup of order $p$. Suppose that $A$ contains no zero-sum
subsequence of length $p$ or $2p$. Then
$\left|\left(\Sigma_p(A)\cup\Sigma_{2p}(A)\right)\cap H\right|
\ge N-2p+2.$
\end{lemma}

\begin{proof} We may assume that $N\ge 2p-1$, since otherwise
$N-2p+2\leq0$, and hence the conclusion follows trivially.

Write $V=H\bigoplus K$ where $K\simeq C_p$. Let $\theta_1$ and $\theta_2$ be the canonical projections of $V$ onto the components $H$ and $K$, respectively.  Then we consider the polynomial ring $\mathbb F_p[X_1,\ldots,X_{N}]$ over the field $\mathbb F_p$ with $N$ indeterminates. Since $C_p$ is isomorphic to the additive group of the field $\mathbb{F}_p$, we can view $\theta_i$ as the map of $V$ to $\mathbb{F}_p$ for $i\in [1,2]$. Then we take the following three polynomials $L(X_1,\ldots,X_N), Q(X_1,\ldots,X_N), J(X_1,\ldots,X_N)\in \mathbb F_p[X_1,\ldots,X_{N}]$, where
$$\begin{aligned}
&L(X_1,\ldots,X_N)=\sum_i\theta_1(a_i)X_i,  \\ 
&Q(X_1,\ldots,X_N)=\sum_i \theta_2(a_i)X_i, \\  
&J(X_1,\ldots,X_N)=\sum_{i=1}^N X_i. \end{aligned}
$$


For $x=(x_1,\ldots,x_N)\in\{0,1\}^N$, let $A_x$ denote the
subsequence of $A$ consisting of all terms $a_i$ whenever $x_i=1$, and we write
$L(x)=L(x_1,\ldots,x_N)$, and similarly for $J(x)$ and $Q(x)$.
Then we define
\begin{equation}\label{equation:Rdefinition}
R=\{L(x):x\in\{0,1\}^N\setminus\{\mathbf 0\},
      \ J(x)=Q(x)=0\}.
\end{equation}

Now we show that
\begin{equation}\label{equation:0notinR}
0\notin R,
\end{equation}
that is $R\subseteq\mathbb F_p\setminus \{0\}$.
Assume to the contrary that $0\in R$. Then there exists
$x=(x_1,\ldots,x_N)\in\{0,1\}^N\setminus\{\mathbf 0\}$ such that
$L(x)=J(x)=Q(x)=0$. Since $J(x)=0$ in $\mathbb F_p$ and $N<3p$,
the number of indices $i\in [1,N]$ with $x_i=1$ is either $p$ or $2p$.
Then $|A_x|\in\{p,2p\}$. Moreover,
$L(x)=Q(x)=0$ implies that
$\sigma(A_x)=0$ in $V=H\oplus K$. Thus $A$ contains a zero-sum
subsequence of length $p$ or $2p$, contradicting the hypothesis.
This proves \eqref{equation:0notinR}.




Then we take the univariate polynomial
\begin{equation}\label{equation:P(Y)definition}
P(Y)=\prod_{r\in R}(Y-r)\in\mathbb F_p[Y],
\end{equation}
where the empty product is understood to be the constant polynomial $1$.
We further take the polynomial $F(X_1,\ldots,X_N) \in \mathbb{F}_p[X_1,\ldots,X_N]$, where \begin{equation}\label{equation:Fdefinition}
F(X_1,\ldots,X_N)=
\left(1-J(X_1,\ldots,X_N)^{p-1}\right)
\left(1-Q(X_1,\ldots,X_N)^{p-1}\right)
P(L(X_1,\ldots,X_N)).
\end{equation}


We claim that
\begin{equation}\label{equation:F(x)=0}
F(x)=0\mbox{ for every } x=(x_1,\ldots,x_N)\in\{0,1\}^N\setminus\{\mathbf 0\}.
\end{equation}
Let $x\in\{0,1\}^N\setminus\{\mathbf 0\}$. If $J(x)\ne0$ or
$Q(x)\ne0$ in $\mathbb F_p$, then, by Fermat's
little theorem, the first or the second factor in
\eqref{equation:Fdefinition} vanishes, respectively. Hence, we may assume that $J(x)=Q(x)=0$. By
\eqref{equation:Rdefinition}, we then have $L(x)\in R$, and therefore
$P(L(x))=0$ by \eqref{equation:P(Y)definition}. Thus the third factor
in \eqref{equation:Fdefinition} vanishes. Consequently, $F(x)=0$. This proves \eqref{equation:F(x)=0}.

On the other hand, since $J(\mathbf 0)=Q(\mathbf 0)=L(\mathbf 0)=0$, it follows from \eqref{equation:0notinR}, \eqref{equation:P(Y)definition} and \eqref{equation:Fdefinition} that $$F(\mathbf 0)=P(L(\mathbf 0))=P(0)\neq 0.$$ Combined with \eqref{equation:F(x)=0}, we derive that $F(X_1,\ldots,X_N)$ and
$P(0)\prod_{i=1}^N(1-X_i)$ have the same value at every point of
$\{0,1\}^N$.

Since $x_i^2=x_i$ for every $x_i\in\{0,1\}$, by taking the remainder modulo $X_i^2-X_i$ for each $i\in[1,N]$,
we obtain a multilinear polynomial
$\widetilde F(X_1,\ldots,X_N)\in
\mathbb F_p[X_1,\ldots,X_N]$
such that
$$\widetilde F(x)=F(x)
\qquad\text{for every }x\in\{0,1\}^N.$$
The polynomial $P(0)\prod_{i=1}^N(1-X_i)$ is multilinear and has the
same values as $\widetilde F$ at every point of $\{0,1\}^N$. Hence,
by Lemma~\ref{lemma:multilinearreduction},
$$\widetilde F(X_1,\ldots,X_N)=P(0)\prod_{i=1}^N(1-X_i).$$
Since $P(0)\ne0$, we have $$\deg\widetilde F=N.$$ By \eqref{equation:P(Y)definition} and the definition of $L(X_1,\ldots,X_N)$, we have that $$\deg P(L(X_1,\ldots,X_N))\leq \deg P(Y)=|R|.$$
Moreover, taking the remainder
modulo $X_i^2-X_i$ cannot increase the total degree. Therefore, by
\eqref{equation:Fdefinition}, we conclude that
$N=\deg\widetilde F\leq\deg F\leq(p-1)+(p-1)+|R|=2p-2+|R|$.
Consequently,
\begin{equation}\label{equation:Rbound}
|R|\ge N-2p+2.
\end{equation}


Finally, we identify the set $R$.  If
$x=(x_1,\ldots,x_N)\in\{0,1\}^N\setminus\{\mathbf 0\}$ satisfies $J(x)=0$, then
$0<|A_x|<3p$ and $|A_x|\equiv0\pmod p$, so
$|A_x|\in\{p,2p\}$. Moreover, $Q(x)=0$ means that
$\sigma(A_x)\in H$, and in this case
$L(x)=\sigma(A_x)$. Hence $R=
\left(\Sigma_p(A)\cup\Sigma_{2p}(A)\right)\cap H.$
Together with \eqref{equation:Rbound}, this gives $\left|
\left(\Sigma_p(A)\cup\Sigma_{2p}(A)\right)\cap H
\right|
\ge N-2p+2,$
as desired.
\end{proof}

\end{document}
